\documentclass[../main.tex]{subfiles}
\begin{document}

\begin{center}
    \Large{\textbf{ABSTRACT}}\\
\end{center}
\vspace{1cm}
Trong đào tạo nội bộ tại doanh nghiệp phần mềm, các module học tập cần được cải tiến liên tục về mục tiêu, đánh giá và hoạt động, chứ không chỉ thiết kế một lần. Các \gls{LLM} đủ mạnh hiện nay đã có thể đảm đương gần như toàn bộ pipeline \gls{ISD}, chứ không chỉ hỗ trợ từng tác vụ riêng lẻ. Khi năng lực của model nền dần bão hoà, khác biệt về chất lượng đầu ra chuyển dịch sang lớp harness bao quanh model, thậm chí lớp này càng có giá trị khi kết hợp với \gls{SLM} nhằm tối ưu chi phí trong thực tế doanh nghiệp. Tuy nhiên, các harness cho ISD hiện nay phần lớn vẫn tối ưu cho chất lượng artifact tại lần sinh đầu tiên, chứ chưa hướng đến khả năng tinh chỉnh hiệu quả theo yêu cầu chỉnh sửa về sau. Vì không duy trì mối quan hệ tường minh giữa các thành phần thiết kế, mỗi lượt cải tiến tiếp theo khó truy vết ảnh hưởng của một thay đổi, nên thường phải rà soát hoặc sinh lại trên diện rộng — một cách làm kém hiệu quả so với nhu cầu cải tiến liên tục đặt ra ở trên.

Luận văn đề xuất AlignmentGraph-ISD, một harness biểu diễn tường minh alignment giữa các artifact thiết kế dưới dạng đồ thị, nhờ đó kiểm tra alignment ở từng liên kết, sửa lỗi và tinh chỉnh có mục tiêu mà không phải sinh lại toàn bộ khoá học. Harness điều phối nhiều agent chuyên biệt (multi-agent) quanh đồ thị này, kết hợp agentic RAG truy xuất tri thức theo domain cho ứng dụng thực tế.

ISD-Agent-Bench được lựa chọn làm benchmark đánh giá chất lượng cho lần sinh đầu tiên, trong đó AlignmentGraph-ISD được tích hợp như một agent mới. Kết quả trên hai lượt benchmark với hai model khác nhau cho thấy hệ thống dẫn đầu hoặc thuộc nhóm dẫn đầu về điểm tổng hợp so với sáu agent so sánh và dẫn đầu điểm ADDIE ở cả hai lượt, dù còn khoảng cách về điểm trajectory so với ReAct-ADDIE. Một case study trong miền IT/Software Engineering được sử dụng để đánh giá chất lượng và khả năng tinh chỉnh đa phiên trên cùng một chương trình đào tạo, thông qua một khung process metric; \textcolor{red}{kết quả đo sẽ được bổ sung sau khi các phiên tinh chỉnh được thực thi}.
\begin{flushright}
\begin{tabular}{@{}c@{}}
Student\\
\textit{(Signature and full name)}
\end{tabular}
\end{flushright}

\end{document}